% TOP_PROBLEM: {"problemId": "problem.spectral-determination-of-convex-planar-domains", "canonicalTitle": "Spectral Determination of Convex Planar Domains", "releaseRank": 492, "primaryDomain": "analysis_pde", "primaryDomainLabel": "Analysis & PDE", "displayStatus": "open_with_solved_subcases", "releaseStatus": "open_with_solved_subcases"}
% !TeX program = lualatex
% Definition and sources only; this file is not an LLM solution attempt.
\documentclass[11pt]{article}
\usepackage[margin=25mm]{geometry}
\usepackage{fontspec}
\setmainfont{Latin Modern Roman}
\usepackage{amsmath,amssymb,mathtools}
\usepackage{unicode-math}
\setmathfont{Latin Modern Math}
\usepackage{xurl,hyperref}
\hypersetup{colorlinks=true,urlcolor=blue}
\setcounter{secnumdepth}{0}
\setlength{\parindent}{0pt}
\setlength{\parskip}{5pt}
\setlength{\emergencystretch}{3em}
\begin{document}
\section*{\texorpdfstring{Spectral Determination of Convex Planar Domains}{Spectral Determination of Convex Planar Domains}}\label{492}
\subsection{Definitions and mathematical statement}
For a bounded open convex domain $\Omega\subset{\mathbb{R}}^2$, the Dirichlet Laplacian is the self-adjoint operator associated with the quadratic form $u\mapsto\int_\Omega|\nabla u|^2$ on $H_0^1(\Omega)$. Write its discrete spectrum, with multiplicities, as
\[
 0<\lambda_1(\Omega)\le\lambda_2(\Omega)\le\cdots.
\]
The question is whether
\[
 \lambda_j(\Omega_1)=\lambda_j(\Omega_2)\ (j\ge1)
 \quad\Longrightarrow\quad \Omega_2=U\Omega_1+b
\]
for an orthogonal matrix $U$ and translation $b$. Reflections are included among congruences. The entire spectrum with multiplicities is required; equality of finitely many eigenvalues or heat invariants is weaker. The convexity assumption must be retained when considering examples of isospectral domains.
\par\smallskip\textit{Formulation and concept references:} \href{https://arxiv.org/html/2406.18369v2}{[S1]}.

\subsection{Short English statement}
Does the full Dirichlet spectrum determine a bounded convex planar domain uniquely up to Euclidean congruence?
\subsection{Sources}
\begin{itemize}
\item[S1]G. Mardby and J. Rowlett, 112 years of listening to Riemannian manifolds, section on planar domains (2024).
\url{https://arxiv.org/html/2406.18369v2}.
\textit{Formulation source.}
\end{itemize}
\end{document}
